\section{Introduction}
\label{sec:intro}

A \emph{closed meander} of order $n$ is a closed planar curve crossing a fixed
line at $2n$ points, up to homeomorphism. Cutting the curve along the line
produces two noncrossing perfect matchings of the crossing points, one drawn
above the line and one below; the curve is connected exactly when the two
matchings form a single cycle. The closed meandric numbers
$M_1,M_2,\dots=1,2,8,42,262,\dots$ (OEIS A005315) and the open meandric numbers
$\Open(m)$ (OEIS A005316) have been studied since Lando and Zvonkin
\cite{LandoZvonkin1992} and the work of Di Francesco, Golinelli and Guitter
\cite{DFGG1997Arch,DFGG1996Direct,DFGG1997TL}. No polynomial-time exact counting
algorithm is known. Classical high-order computations use transfer-matrix
methods \cite{Jensen2000,DFGJ2000,Jensen2025}. Jensen reports empirical
signature growth near $2.5^{n}$ in the order $n$ \cite{Jensen2025}; this is
a fitted observation rather than a proved complexity bound.

This paper proves the following.

\begin{theorem}[Main theorem]
\label{thm:main}
\lean{Meanders.FirstCrossing.evaluateJoint_correct, Meanders.FirstCrossing.count_eq, Meanders.FirstCrossing.carrierCard_le}
\uses{thm:cost, cor:balance, thm:joint, cor:open}
There is a deterministic algorithm that, on input $n\ge1$, computes the exact
integers $M_n=\Closed(n)$, $\Open(2n-1)$ and $\Open(2n)$ using
$\Ostar(2^{n})$ bit operations and $\Ostar(2^{n})$ bits of space.
More precisely, for every threshold $1\le K\le n$ the algorithm runs in
$\operatorname{poly}(n)\bigl(4^{K-1}+4^{n-K}\bigr)$ bit time and space, and the
choice $K=\lfloor n/2\rfloor+1$ gives the stated bound. Consequently
$\Open(m)$ can be computed in $\Ostar(2^{\lceil m/2\rceil})=\Ostar(\sqrt2^{\,m})$
for every $m$.
\end{theorem}

The construction combines well-known ingredients: Dyck encodings of noncrossing
matchings, planar connectivity transfer with local cup/cap surgery, remaining-budget
feasibility, weighted state merging, and the classical closed/open counting
identities. The contribution is the way these are made compatible:

\begin{enumerate}[nosep,label=(\roman*)]
\item a \emph{disjoint source partition} by the first cut at which the combined
  original Dyck height $k_j=(h_P(j)+h_Q(j))/2$ reaches $K$, which leaves a
  bounded-height complement (\LOW) and $O(n^2)$ first-hit sectors (\HIGH);
\item in each \HIGH\ sector, an \emph{asymmetric inward scan} of both halves of
  the cut whose exact down budgets make an oldest prefix of the pending arches
  \emph{forced}, so that guaranteed through arches are joined as soon as both
  ends are known;
\item a \emph{four-seam carrier}, a single cyclic noncrossing pairing with four
  groups whose separators are determined by four counters, together with a proof
  that the retained boundary never exceeds $2(n-K)+O(1)$ endpoints at any stage;
\item a \emph{shared observer} for the even open count, based on the identity
  $e(Q)=r_L+r_R+\ind{h_Q(L)>0}$ for the number of exterior arches of a matching
  cut at any position $L$.
\end{enumerate}

Item (iii) is the technical center: a bounded-height cap alone leaves the
complementary pairs uncounted, and retaining two independent halves would destroy
the exponent. We do not claim priority for the individual transfer operations or
for the counting identities, which are due to earlier authors
(\cref{sec:related}).

The algorithm has been implemented in Rust. The reference evaluator's
correctness, the width bounds of \cref{sec:width} and a carrier-cardinality
bound with the exponential factors $4^{K-1}$ and $4^{n-K}$ are formally
verified in Lean~4; the polynomial-factor bit-time and bit-space accounting
of \cref{thm:main} is proved on paper in \cref{sec:width}, and the optimized
Rust implementation is validated differentially (\cref{sec:evidence}). The
implementation reproduces the published values through $M_{28}$ and
$\Open(55)$; a production sweep at ranks $27$ to $30$ also returned
$M_{29}$, $M_{30}$ and $\Open(56)$ to $\Open(60)$ as production outputs
without independent check (\cref{sec:evidence}).

\paragraph{Organization.}
\Cref{sec:prelim} fixes notation and the closed/open conventions.
\Cref{sec:partition} defines the partition, \cref{sec:cut} the cut bijection and
the forced-prefix lemma, \cref{sec:carrier} the carrier and the transfer,
\cref{sec:width} the width and cost bounds, and \cref{sec:open} the open
observer. \Cref{sec:evidence} summarizes formalization, implementation and
measurements; \cref{sec:related} discusses related work. The appendices contain
the port-bound computations, the Lean declaration map and the artifact pointer.
